
% 附录源 PDF 第 18—22 页；A.4 导言和 A.4.1。
\section{在线靶向学习}

本节将推导形如
\begin{equation}
\begin{aligned}
\Psi_J(P^{N*})-\Psi_J(P_0^N)
={}&-\frac1N\sum_{i=1}^NP_{0,o(i)}\bar D(P^{N*})\\
&+R_{2,N}(\theta_N^*,\theta_0)
\end{aligned}
\tag{A.11}
\end{equation}
的一阶展开，其中 $P_{0,o(i)}$ 表示给定 $\bar O(i-1)$ 时关于 $O(i)$ 的条件期望，
$R_{2,N}(\theta_N^*,\theta_0)$ 为二阶余项。定理 A.4 给出了这一结果。
为建立该一阶展开，我们首先证明恒等式
\[
\begin{aligned}
\Psi_J(P^{N*})-\Psi_J(P_0^N)
={}&-\frac1N\sum_{i=1}^NP_0^N\bar D(P^{N*})\\
&+R_{21}(\theta_N^*,\theta_0),
\end{aligned}
\]
其中明确给出二阶余项 $R_{21}(\theta_N^*,\theta_0)$。定理 A.3 陈述这一结果。
注意，这里涉及有效影响曲线的边际期望，而所需的是条件期望。随后在第二小节研究额外的余项
\[
R_{22,N}(\theta_N^*,\theta_0)
 =\frac1N\sum_{i=1}^N(P_{0,o(i)}-P_0^N)
 \{\bar D(\theta_N^*)-\bar D(\theta_0)\},
\]
其中 $R_{2,N}()=R_{21}()+R_{22,N}()$。接着说明，在时间序列依赖结构满足一定条件时，$R_{22,N}$ 也是二阶余项。

恒等式 (A.11) 是分析 TMLE 的基础。具体而言，把式 (A.11) 与解方程
$0=N^{-1}\sum_{i=1}^N\bar D(\theta_N^*)=0$ 的 TMLE $\theta_N^*$ 结合，就得到一阶展开
\[
\begin{aligned}
\Psi_J(P^{N*})-\Psi_J(P_0^N)
={}&\frac1N\sum_{i=1}^N\{\bar D(P^{N*})-P_{0,o(i)}\bar D(P^{N*})\}\\
&+R_{2,N}(\theta_N^*,\theta_0).
\end{aligned}
\]
右侧主导项是以 $\theta_N^*$ 为指标的鞅过程，因此可在较弱条件下证明其收敛到正态极限分布。
随后，我们为在线估计建立类似的一阶展开：
\[
\begin{aligned}
-\frac1N\sum_{i=1}^NP_{0,o(i)}\bar D_{g^*}(\theta_{i-1})
={}&\frac1N\sum_{i=1}^N\Psi_{g^*}(\theta_{i-1})-\Psi_{g^*}(\theta_0)\\
&+\bar R_{21,g^*,N}+\bar R_{22,g^*,N}.
\end{aligned}
\]
这一结果同样分两步推导：首先基于边际期望 $P_0^N\bar D_{g^*}(\theta_{i-1})$ 建立只涉及余项 $\bar R_{21,g^*,N}$ 的结果；然后研究余项 $\bar R_{22,g^*,N}$，并指出，要控制它还需增加一个平稳性假设。

\subsection{基于有效影响曲线边际期望的目标参数一阶展开}

注意，
\[
P_{g^*,o(t)}f=P_{w(t)}P_{g^*,a(t)}P_{y(t)}f.
\]
为简化记号，也将 $P_{g^*}$ 记为 $P^*$。有
\[
\Psi_{g^*}(P)=\left(\prod_{t=1}^{\tau}P_{g^*,o(t)}\right)Y(\tau).
\]

\begin{bookresult}{定理}{A.3}
令
\[
\begin{aligned}
\Delta_w(P^*_{o(s)},P^*_{0,o(s)})
 &=(P_{w(s)}-P_{0,w(s)})P^*_{0,a(s)}P_{0,y(s)},\\
\Delta_a(P^*_{o(s)},P^*_{0,o(s)})
 &=P_{w(s)}(P^*_{a(s)}-P^*_{0,a(s)})P_{0,y(s)},\\
\Delta_y(P^*_{o(s)},P^*_{0,o(s)})
 &=P_{w(s)}P_{a(s)}(P_{y(s)}-P_{0,y(s)}),\\
\Phi_{x(s)}(x,c_x)
 &=E(Y_{g^*}(\tau)\mid X(s)=x,C_x(s)=c_x).
\end{aligned}
\]
则
\[
-P_0^ND_{g^*}^N(P)=\Psi_{g^*}(P)-\Psi_{g^*}(P_0)+R_{21,g^*}(P,P_0),
\]
其中 $R_{21,g^*}(P,P_0)=R_{21a,g^*}(P,P_0)+R_{21b,g^*}(P,P_0)$，且
\begin{equation}
\begin{aligned}
R_{21a,g^*}(\theta,\theta_0)
={}&\sum_{x\in\{w,a,y\}}\sum_{s=1}^{\tau}
 \left\{\prod_{t=1}^{s-1}P^*_{o(t)}\right\}
 \Delta_x(P^*_{o(s)},P^*_{0,o(s)})\\
&\quad\times\Biggl\{\sum_{j=s+1}^{\tau}
 \prod_{l=s+1}^{j-1}P^*_{o(l)}(P^*_{o(j)}-P^*_{0,o(j)})\\
&\hspace{12em}\times\prod_{l=j+1}^{\tau}P^*_{0,o(l)}\Biggr\}Y(\tau),
\end{aligned}
\tag{A.12}
\end{equation}
\[
R_{21b,g^*}(P,P_0)
 =\sum_{x\in\{w,a,y\}}\sum_{s=1}^{\tau}R_{21b,x(s),g^*}(P,P_0),
\]
\[
\begin{aligned}
&R_{21b,x(s),g^*}(P,P_0)\\
&=\int_c\int_x\left(\frac{\bar h_{c_x,0}}{\bar h_{c_x}}(c)-1\right)
 \Phi_{x(s)}(x,c)\\
&\qquad\times\mathrm d(P^*_{x(s)}-P^*_{0,x(s)})(x\mid c)
 h^*_{c_x(s)}(c)\,\mathrm d\mu(c).
\end{aligned}
\]
因此，回顾
$\Psi_J(P)=J^{-1}\sum_{j=1}^J\Psi_{g_j^*}(P)$ 和
$D^N(P)=J^{-1}\sum_{j=1}^JD^N_{g_j^*}(P)$，有
\[
-P_0^ND^N(P)=\Psi_J(P)-\Psi_J(P_0)+R_{21}(\theta,\theta_0),
\]
其中 $R_{21}(\theta,\theta_0)=J^{-1}\sum_{j=1}^JR_{21,g_j^*}(\theta,\theta_0)$。
\end{bookresult}

\begin{proof}
定义
\[
\begin{aligned}
&\Psi_{g^*,P,P_0,w(s)}(P^*_{1,w(s)})\\
&\equiv\left\{\prod_{t=1}^{s-1}P^*_{o(t)}\right\}
 P_{1,w(s)}P^*_{0,a(s)}P_{0,y(s)}
 \left\{\prod_{t=s+1}^{\tau}P^*_{0,o(t)}\right\}Y(\tau),\\
&\Psi_{g^*,P,P_0,a(s)}(P^*_{1,a(s)})\\
&\equiv\left\{\prod_{t=1}^{s-1}P^*_{o(t)}\right\}
 P_{w(s)}P^*_{1,a(s)}P_{0,y(s)}
 \left\{\prod_{t=s+1}^{\tau}P^*_{0,o(t)}\right\}Y(\tau),\\
&\Psi_{g^*,P,P_0,y(s)}(P^*_{1,y(s)})\\
&\equiv\left\{\prod_{t=1}^{s-1}P^*_{o(t)}\right\}
 P_{w(s)}P^*_{a(s)}P_{1,y(s)}
 \left\{\prod_{t=s+1}^{\tau}P^*_{0,o(t)}\right\}Y(\tau).
\end{aligned}
\]
注意，
\[
\begin{aligned}
&\Psi_{g^*}(P)-\Psi_{g^*}(P_0)\\
&=\sum_{x\in\{w,a,y\}}\sum_{s=1}^{\tau}
 \Biggl(\left\{\prod_{t=1}^{s-1}P^*_{o(t)}\right\}
 \Delta_x(P^*_{o(s)},P^*_{0,o(s)})\\
&\hspace{13em}\times\left\{\prod_{t=s+1}^{\tau}P^*_{0,o(t)}\right\}\Biggr)Y(\tau)\\
&=\sum_{x\in\{w,a,y\}}\sum_{s=1}^{\tau}
 \bigl\{\Psi_{g^*,P,P_0,x(s)}(P^*_{x(s)})
 -\Psi_{g^*,P,P_0,x(s)}(P^*_{0,x(s)})\bigr\}.
\end{aligned}
\]
定义 $\Psi_{g^*,P,x(s)}=\Psi_{g^*,P,P,x(s)}$，并定义二阶项
\[
\begin{aligned}
&-R_{21a,g^*}(\theta,\theta_0)\\
&=\sum_{x\in\{w,a,y\}}\sum_{s=1}^{\tau}
 \Biggl(\left\{\prod_{t=1}^{s-1}P^*_{o(t)}\right\}
 \Delta_x(P^*_{o(s)},P^*_{0,o(s)})
 \left\{\prod_{t=s+1}^{\tau}P^*_{0,o(t)}\right\}\Biggr)Y(\tau)\\
&\quad-\sum_{x\in\{w,a,y\}}\sum_{s=1}^{\tau}
 \Biggl(\left\{\prod_{t=1}^{s-1}P^*_{o(t)}\right\}
 \Delta_x(P^*_{o(s)},P^*_{0,o(s)})
 \left\{\prod_{t=s+1}^{\tau}P^*_{o(t)}\right\}\Biggr)Y(\tau)\\
&=-\sum_{x\in\{w,a,y\}}\sum_{s=1}^{\tau}
 \left\{\prod_{t=1}^{s-1}P^*_{o(t)}\right\}
 \Delta_x(P^*_{o(s)},P^*_{0,o(s)})\\
&\qquad\times\Biggl\{\sum_{j=s+1}^{\tau}
 \prod_{l=s+1}^{j-1}P^*_{o(l)}(P^*_{o(j)}-P^*_{0,o(j)})
 \prod_{l=j+1}^{\tau}P^*_{0,o(l)}\Biggr\}Y(\tau).
\end{aligned}
\]
于是，
\[
\begin{aligned}
&\Psi_{g^*}(P)-\Psi_{g^*}(P_0)\\
&=\sum_{x\in\{w,a,y\}}\sum_{s=1}^{\tau}
 \bigl(\Psi_{g^*,P,x(s)}(P^*_{x(s)})
 -\Psi_{g^*,P,x(s)}(P^*_{0,x(s)})\bigr)
 +R_{21a,g^*}(\theta,\theta_0).
\end{aligned}
\]
另外，
\[
D_{g^*}^N(P)(O^N)
 =\sum_{x\in\{y,a,w\}}\sum_{s=1}^{\tau}\frac1N\sum_{i=1}^N
 \bar D_{x(s)}(P)(X(i),C_x(i)),
\]
其中 $\bar D_{x(s)}$ 是有效影响曲线中对应于给定 $C_x(s)$ 时 $X(s)$ 条件分布得分的分量。
令
\[
D_{x(s)}^N(P)=\frac1N\sum_{i=1}^N\bar D_{x(s)}(P)(X(i),C_x(i)).
\]
我们将证明
\begin{equation}
\begin{aligned}
-P_0^ND_{x(s)}^N(P)
={}&\Psi_{g^*,P,x(s)}(P^*_{x(s)})
 -\Psi_{g^*,P,x(s)}(P^*_{0,x(s)})\\
&+R_{21b,x(s),g^*}(P,P_0).
\end{aligned}
\tag{A.13}
\end{equation}
这就给出以下推导：
\[
\begin{aligned}
&-P_0^ND_{g^*}^N(P)\\
&=\sum_{x\in\{y,a,w\}}\sum_{s=1}^{\tau}-P_0^ND_{x(s)}^N(P)\\
&=\sum_{x\in\{y,a,w\}}\sum_{s=1}^{\tau}
 \Bigl\{\Psi_{g^*,P,x(s)}(P^*_{x(s)})
 -\Psi_{g^*,P,x(s)}(P^*_{0,x(s)})
 +R_{21b,x(s),g^*}(P,P_0)\Bigr\}\\
&=\sum_{x\in\{y,a,w\}}\sum_{s=1}^{\tau}
 \bigl\{\Psi_{g^*,P,x(s)}(P^*_{x(s)})
 -\Psi_{g^*,P,x(s)}(P^*_{0,x(s)})\bigr\}\\
&\quad+\sum_{x\in\{y,a,w\}}\sum_{s=1}^{\tau}R_{21b,x(s),g^*}(P,P_0)\\
&=\Psi_{g^*}(P)-\Psi_{g^*}(P_0)-R_{21a,g^*}(\theta,\theta_0)\\
&\quad+\sum_{x\in\{y,a,w\}}\sum_{s=1}^{\tau}R_{21b,x(s),g^*}(P,P_0)\\
&\equiv\Psi_{g^*}(P)-\Psi_{g^*}(P_0)-R_{21a,g^*}(\theta,\theta_0)
 +R_{21b,g^*}(P,P_0),
\end{aligned}
\]
从而证明了定理。因此，只需再建立式 (A.13)。有
\[
\begin{aligned}
&-P_0^ND_{x(s)}^N(P)\\
&=-P_0^N\frac1N\sum_{i=1}^N
 \frac{h^*_{c_x(s)}(C_x(i))}{\bar h_{c_x}(C_x(i))}\\
&\quad\times\Bigl\{E(Y_{g^*}(\tau)\mid X(s)=X(i),C_x(s)=C_x(i))\\
&\hspace{7em}-E(Y_{g^*}(\tau)\mid C_X(s)=C_x(i))\Bigr\}\\
&=P_0^N\frac1N\sum_{i=1}^N
 \frac{h^*_{c_x(s)}(C_x(i))}{\bar h_{c_x}(C_x(i))}
 \int_x\Phi_{x(s)}(x,C_x(i))\\
&\hspace{8em}\times\mathrm d(P^*_{x(s)}-P^*_{0,x(s)})(x\mid C_x(i))\\
&=\frac1N\sum_{i=1}^N\int_c
 \frac{h^*_{c_x(s)}(c)}{\bar h_{c_x}(c)}
 \int_x\Phi_{x(s)}(x,c)\\
&\hspace{6em}\times\mathrm d(P^*_{x(s)}-P^*_{0,x(s)})(x\mid c)
 h_{c_x(i),0}(c)\,\mathrm d\mu(c)\\
&=\int_c\int_x\frac{\bar h_{c_x,0}(c)}{\bar h_{c_x}(c)}
 \Phi_{x(s)}(x,c)\,\mathrm d(P^*_{x(s)}-P^*_{0,x(s)})(x\mid c)
 h^*_{c_x(s)}(c)\,\mathrm d\mu(c)\\
&=\int_c\int_x\left(\frac{\bar h_{c_x,0}}{\bar h_{c_x}}(c)-1\right)
 \Phi_{x(s)}(x,c)\,\mathrm d(P^*_{x(s)}-P^*_{0,x(s)})(x\mid c)
 h^*_{c_x(s)}(c)\,\mathrm d\mu(c)\\
&\quad+\int_{x,c}\Phi_{x(s)}(x,c)
 \,\mathrm d(P^*_{x(s)}-P^*_{0,x(s)})(x\mid c)
 h^*_{c_x(s)}(c)\,\mathrm d\mu(c)\\
&\equiv R_{21b,x(s),g^*}(P,P_0)
 +\Psi_{g^*,P,x(s)}(P^*_{x(s)})
 -\Psi_{g^*,P,x(s)}(P^*_{0,x(s)}).
\end{aligned}
\]
证明完毕。
\end{proof}
